\section{Domains and conditions of the realizations}
\label{v:integracion}
Algebraic exclusion, its transport under refinement, and its statistical
applications are related results with distinct domains. The general
annihilation identity \eqref{v:fock:eq:transporte-aniquilacion} makes
the adjoint of the transport explicit; its isometric specialization
retains the necessary condition on the map between the one-particle
spaces. Equilibrium occupations additionally use a grading and a
Hamiltonian compatible with the trajectory representation.

The finite dynamical realization is developed in
Section~\ref{sec:hmtf2-cristal-temporal}. The area operator, modular
Hamiltonian, and sectoral purification are developed in
Section~\ref{v:ant:sec:area} and its continuations. Sections
\ref{v:ant:sec:action} and \ref{v:ant:sec:ii-kb-aut-desarrollo}
contain the constructions of action and the Boltzmann--temperature
pair used by the radiative realization. The variational characterization
of the history measure and convergence of the spectral sums link,
respectively, incidence to entropy and discrete modes to radiation.
The radiation section presents its cavity, dispersion, polarization,
and equilibrium hypotheses, together with the spectral proofs it uses.
Section~\ref{v:sec:enlace-velocidad-constitutiva} also establishes
the identity of the kinematic and constitutive realizations of speed,
its dimensional normalization, and its preservation under amplification
of the two-sheet space.

Each composition retains its stated mathematical hypotheses and physical
realization. In particular, the harmonic modes of the cellular complex,
connections modulo gauge, and radiative modes retain the domains
distinguished in Section~\ref{v:rec:completaciones:section:08};
the radiative realization explicitly uses its dispersion relation and
polarizations.

Subsection~\ref{v:subsec:limite-termodinamico} proves, for the three
observables \(F_N,F_E,F_Z\), the passage from three-dimensional periodic
sums to radial integrals as \(L\to\infty\). Its domain requires \(a,b>0\),
retains the fixed background constant mode, and uniformly controls the
origin and the spectral tail. The result yields the number, excitation-energy,
and logarithmic partition-function densities of the transverse modes; it is
distinct from the finite spectral-capacity block.

Section~\ref{v:sec:capacidad-espectral-termica} has as its domain the finite
graded carrier, the microscopic and aggregated trees, the reversible memory
lift, the capacity clock, and the thermometric charts declared there. Its
spectral, transport, and variational identities are proved within those
domains; physical recognition enters afterwards. The typed dimensional
chart preserves the provenance of the marked intensity, energy chart, and
physical temperature.

Section~\ref{v:n19:sec:incertidumbre-fourier} works on the finite
realization \(\ell^2((\mathbb Z/27\mathbb Z)^2)\) and its unitary Fourier
transform. Its entropies are those of the distributions in two conjugate
bases; the dodecaphase readers have the distinct domain
\(\{-1,0,+1\}^{12}\) and provide a complementary ordered-memory result.
Section~\ref{sec:f051-prime-vacancy-observer} establishes the exact modal
identity for each finite \(X\); its historical witness reaches \(10^8\),
and the documented independent reproduction reaches \(10^6\). Promotion
to the critical bound requires estimates uniform in the modes.
